% 学术内容检测模块整体业务流程（对应 03-5-content-detection.tex 全部功能点 FR-TPJC-0001～0003、FR-LWJC-0001～0005、FR-PLJC-0001～0002、FR-NRJC-0001～0002）
\begin{tikzpicture}[
  box/.style={draw,rounded corners,align=center,minimum width=1.95cm,inner sep=1.4pt,font=\tiny},
  arr/.style={-{Stealth[length=1.4mm]},semithick}
]
  \node[box] (L) at (0,0) {登录/任务就绪};
  \node[box] (S) at (0,-1.1) {检测任务进入\\学术内容检测模块};

  \node[box] (T1) at (-4.6,-2.45) {提取图片\\\textbf{TPJC-0001}};
  \node[box] (P1) at (0,-2.45) {全篇论文AIGC检测\\\textbf{LWJC-0001}};
  \node[box] (R1) at (4.6,-2.45) {Review检测\\\textbf{PLJC-0001}};

  \node[box] (T2) at (-4.6,-3.75) {学术图像AI鉴伪\\\textbf{TPJC-0002}};
  \node[box] (P2) at (0,-3.75) {AI生成段落识别\\\textbf{LWJC-0002}};
  \node[box] (R2) at (4.6,-3.75) {模板化评审倾向\\识别 \textbf{PLJC-0002}};

  \node[box] (T3) at (-4.6,-5.05) {图像可疑区域\\标注 \textbf{TPJC-0003}};
  \node[box] (P3) at (0,-5.05) {段落级AI生成\\概率 \textbf{LWJC-0003}};

  \node[box] (P4) at (0,-6.35) {事实性鉴伪\\子结论 \textbf{LWJC-0004}};
  \node[box] (P5) at (0,-7.65) {事实性鉴伪\\原因说明 \textbf{LWJC-0005}};

  \node[box] (N1) at (0,-9.05) {综合鉴伪结果生成\\\textbf{NRJC-0001}};
  \node[box] (N2) at (0,-10.4) {检测结果可视化解释\\\textbf{NRJC-0002}};

  \node[box] (N3) at (0,-11.75) {多模态联合检测\\\textbf{XNJC-0006}};
  \node[box] (N4) at (0,-13.1) {综合结果融合分析\\\textbf{XNJC-0007}};


  \draw[arr] (L) -- (S);
  \draw[arr] (S) -- (T1);
  \draw[arr] (S) -- (P1);
  \draw[arr] (S) -- (R1);
  \draw[arr] (T1) -- (T2) -- (T3);
  \draw[arr] (P1) -- (P2) -- (P3) -- (P4) -- (P5);
  \draw[arr] (R1) -- (R2);
  \draw[arr] (T3.south) -- ++(0,-0.18) -| (N1.west);
  \draw[arr] (P5) -- (N1);
  \draw[arr] (R2.south) -- ++(0,-0.18) -| (N1.east);
  \draw[arr] (N1) -- (N2);
  \draw[arr] (N2) -- (N3);
  \draw[arr] (N3) -- (N4);
\end{tikzpicture}
